%% Le glossaire est un élément optionnel qui peut être utile pour expliquer des termes techniques.

\newglossaryentry{heig-vd}{
    name=HEIG-VD,
    description={Haute École d'Ingénierie et de Gestion du canton de Vaud}
}

\newglossaryentry{hes-so}{
    name=HES-SO,
    description={Haute École Supérieure de Suisse Occidentale}
}

\newglossaryentry{latex}{
    name=latex,
    description={Un langage et un système de composition de documents}
}

\newglossaryentry{maths}{
    name=mathematics,
    description={Les mathematiques sont ce que les mathématiciens fonts}
}

\newglossaryentry{DAC}{
    name=DAC,
    description={Convertisseur numérique-analogique}
}

\newglossaryentry{i2s}{
    name={I\textsuperscript{2}S},
    description={Inter-IC Sound. Bus numérique série dédié au transport d'audio entre circuits intégrés. Il transmet les échantillons sur trois signaux, une horloge de bit, une horloge de mot et une ligne de données, sans protocole d'adressage}
}

\newglossaryentry{bck}{
    name={BCK},
    description={Bit Clock. Horloge du bus I\textsuperscript{2}S qui cadence la transmission de chaque bit. Sa fréquence vaut le produit de la fréquence d'échantillonnage, du nombre de canaux et de la largeur d'un slot}
}

\newglossaryentry{lrck}{
    name={LRCK},
    description={Left Right Clock. Horloge du bus I\textsuperscript{2}S qui désigne le canal auquel appartient le bit courant, l'état bas correspondant au canal gauche et l'état haut au canal droit. Elle bat une fois par échantillon stéréo, donc à la fréquence d'échantillonnage}
}

\newglossaryentry{mclk}{
    name={MCLK},
    description={Master Clock. Horloge haute fréquence fournie optionnellement à un convertisseur pour cadencer son électronique interne. Lorsqu'elle n'est pas câblée, le convertisseur doit la reconstruire à partir de l'horloge de bit au moyen d'une boucle à verrouillage de phase}
}

\newglossaryentry{trame}{
    name={trame audio},
    description={Ensemble des échantillons correspondant à un même instant sur tous les canaux. Une trame stéréo dont chaque canal occupe 32 bits pèse ainsi 8 octets. Le débit du bus se mesure en trames par seconde, valeur égale à la fréquence d'échantillonnage}
}

\newglossaryentry{slot}{
    name={slot},
    description={Intervalle de temps réservé à un canal dans une trame du bus I\textsuperscript{2}S. Sa largeur, exprimée en bits, est indépendante de la résolution réelle des échantillons transmis, les bits inutilisés étant remplis de zéros}
}

\newglossaryentry{msbjust}{
    name={justification MSB},
    description={Placement d'un échantillon en tête de slot, bit de poids fort en premier, les bits de poids faible restants étant complétés par des zéros. Cette convention permet de transmettre un échantillon de résolution quelconque dans un slot de largeur fixe}
}

\newglossaryentry{pll}{
    name={PLL},
    description={Phase-Locked Loop, ou boucle à verrouillage de phase. Circuit asservi qui synthétise une horloge de sortie dont la fréquence est un multiple rationnel de celle d'une horloge de référence}
}

\newglossaryentry{apll}{
    name={APLL},
    description={Audio PLL. Seconde boucle à verrouillage de phase de l'ESP32, dont la fréquence de sortie est finement programmable. Elle permet de synthétiser exactement les horloges des familles de fréquences audio, que la boucle principale du micro-contrôleur ne peut atteindre par division entière}
}

\newglossaryentry{dma}{
    name={DMA},
    description={Direct Memory Access. Mécanisme permettant à un périphérique de lire ou d'écrire en mémoire sans intervention du processeur. Les zones concernées sont décrites par des descripteurs chaînés, chacun désignant un tampon de taille bornée}
}

\newglossaryentry{sink}{
    name={sink audio},
    description={Interface abstraite représentant une destination de flux audio dans le firmware. Elle expose trois opérations, l'ouverture dans un format donné, l'écriture d'un bloc d'échantillons et la fermeture, ce qui rend le producteur d'échantillons indépendant de la sortie réellement utilisée}
}
\newglossaryentry{sd}{
    name={SD},
    description={Secure Digital. Format de carte mémoire amovible de grande capacité. Le lecteur y accède en mode SPI, dont le câblage est plus simple que celui du mode natif au prix d'un débit moindre, sans conséquence pour la lecture audio}
}

\newglossaryentry{psram}{
    name={PSRAM},
    description={Pseudo-Static RAM. Mémoire vive de grande capacité intégrée au module mais externe au micro-contrôleur. Elle accueille les données volumineuses du firmware, à l'exception des tampons manipulés par le contrôleur d'accès direct à la mémoire, qui ne peut pas l'adresser}
}

\newglossaryentry{spi}{
    name={SPI},
    description={Serial Peripheral Interface. Bus série synchrone reliant un contrôleur maître à un ou plusieurs périphériques, chacun sélectionné par une ligne dédiée. Le micro-contrôleur en expose plusieurs instances matérielles indépendantes}
}

\newglossaryentry{pcm}{
    name={PCM},
    description={Pulse-Code Modulation. Représentation numérique directe d'un signal audio par une suite d'échantillons d'amplitude prélevés à intervalle régulier, sans compression. C'est sous ce format que circulent les échantillons entre les décodeurs et le sink}
}

\newglossaryentry{mp3}{
    name={MP3},
    description={MPEG-1 Audio Layer III. Format de compression audio avec pertes, à débit constant ou variable, décodé par le firmware vers la représentation PCM interne}
}

\newglossaryentry{wav}{
    name={WAV},
    description={Waveform Audio File Format. Format de fichier encapsulant des échantillons PCM non compressés, précédés d'un en-tête qui déclare leur fréquence d'échantillonnage, leur résolution et leur nombre de canaux}
}

\newglossaryentry{fatfs}{
    name={FATFS},
    description={Implémentation du système de fichiers FAT destinée aux systèmes embarqués et intégrée à ESP-IDF, par laquelle le firmware accède aux fichiers de la carte SD}
}

\newglossaryentry{a2dp}{
    name={A2DP},
    description={Advanced Audio Distribution Profile. Profil Bluetooth qui régit le transport d'un flux audio stéréo entre une source, qui émet le flux, et un puits, qui le restitue. Sur la radio classique de ce micro-contrôleur, il impose l'emploi du codec SBC}
}

\newglossaryentry{avrcp}{
    name={AVRCP},
    description={Audio/Video Remote Control Profile. Profil Bluetooth de télécommande associé à l'A2DP, par lequel une source peut notamment fixer le volume absolu du récepteur}
}

\newglossaryentry{sbc}{
    name={SBC},
    description={Low Complexity Subband Coding. Codec audio à pertes de faible complexité, seul codec dont la prise en charge est obligatoire dans le profil A2DP. Il opère sur des échantillons de 16 bits}
}

\newglossaryentry{nvs}{
    name={NVS},
    description={Non-Volatile Storage. Service d'ESP-IDF qui conserve de petites données clé-valeur en mémoire flash au travers des cycles d'alimentation, employé par le firmware pour ses préférences persistantes}
}